\section{A quadratic Jensen theorem for a nonabelian Artin factor}
\label{sec:eta-jensen}
Let $F=\mathbb Q(\alpha)$ with $\alpha^3-\alpha-1=0$. The polynomial is
irreducible over $\mathbb Q$ and has discriminant $-23$. Its splitting
field has group $S_3$. Write $\rho$ for the two-dimensional summand of its
three-letter permutation representation. The classical weight-one identity is
\[
 L(s,\rho)=\frac{\zeta_F(s)}{\zeta(s)}=L(s,f),
 \qquad f(z)=\eta(z)\eta(23z).
\]
Here $L(s,f)$ uses the Fourier coefficients of the normalized cusp form.
See van der Blij's result in D. Zagier, \emph{Elliptic Modular Forms and
Their Applications}, Proposition 14 and the discussion following it,
\href{https://www.its.caltech.edu/~matilde/Zagier123ModularForms.pdf}{pp.~43--44}.
In particular the trace at an unramified prime is the number of roots of
$X^3-X-1$ in $\mathbb F_p$ minus one.
Put
\[
 \Lambda_f(s)=\left(\frac{\sqrt{23}}{2\pi}\right)^s\Gamma(s)L(s,f),
 \qquad
 \Lambda_f(1/2+z)=\sum_{n\ge0}\frac{\gamma_n}{n!}z^{2n}.
\]
\textbf{Theorem.} For every $n\ge0$, the polynomial
$\gamma_n+2\gamma_{n+1}X+\gamma_{n+2}X^2$ has two distinct negative
real roots. The conclusion holds independently of Dirichlet GRH.

\textit{Proof.}
After the change of variables $y=e^u/\sqrt{23}$ in the Mellin integral
of $f(iy)$, we obtain
\[
 \Lambda_f(1/2+z)=\int_{\mathbb R}e^{zu}K(u)\,du,
 \qquad K(u)=e^{u/2}\eta(ie^u/\sqrt{23})\eta(i\sqrt{23}e^u).
\]
The integral agrees with the completed Dirichlet series in its initial
half-plane, then gives its entire continuation by the decay below.
Write $h(v)=v/4+\log\eta(ie^v)$ and $a=\frac12\log23$. Then
$\log K(u)=h(u-a)+h(u+a)$. Eta inversion shows that $h$ is even.
For $v\ge0$, the product for $\eta$ gives, with $Y=2\pi e^v$,
\[
 h(v)=\frac v4-\frac Y{24}
       -\sum_{m\ge1}\sigma_{-1}(m)e^{-mY},\qquad
 \sigma_{-1}(m)=\sum_{d\mid m}d^{-1}.
\]
The series and its first four derivatives converge locally uniformly.
The fourth derivative in $v$ of $e^{-mY}$ equals
$e^{-q}q((q-6)q^2+7q-1)$ with $q=mY\ge 2\pi>6$. Consequently
$h^{(4)}(v)<0$ for all real $v$, using evenness for $v<0$.

More generally, let $g(u)=\sum_j r_j[h(u-a_j)+h(u+a_j)]$ for any
nonempty finite list of $a_j\in\mathbb R$, $r_j>0$, and put $K=e^g$.
Then $g$ is even and $g^{(4)}<0$. Thus $g'''(u)<0$ for $u>0$ and
$(u g''-g')'=u g'''<0$. Since $g'(0)=0$,
\[
 \left(\log K(\sqrt{x})\right)''
 =\frac{u g''(u)-g'(u)}{4u^3}<0,\qquad u=\sqrt x>0 .
\]
At both infinities $K$ decays at least as fast as $\exp(-c e^{|u|})$
for some $c>0$, by the same product and symmetry. The preceding
calculation follows the eta-kernel curvature argument of J. B. Conrey
and A. Ghosh, \emph{Tur\'an inequalities and zeros of Dirichlet series
associated with certain cusp forms}, Trans. AMS 342 (1994), 407--419,
\href{https://aimath.org/~kaur/publications/29.pdf}{Secs.~4--7}.
The symmetry of the translations preserves the fourth-derivative sign
for unequal eta dilations.

For completeness, write $V(x)=-\log K(\sqrt{x})$ and
$M(p)=\int_0^\infty x^{p-1}e^{-V(x)}\,dx$ for $p>0$. Under probability
density $x^{p-1}e^{-V(x)}/M(p)$ denote expectation by $E_p$.
Integration by parts gives
$E_{p+1}[V']=p M(p)/M(p+1)$. Since $V'$ is strictly increasing,
\[
 E_{p+2}[V']-E_{p+1}[V']
 =\frac{\operatorname{Cov}_{p+1}(x,V'(x))}{E_{p+1}[x]}>0.
\]
The last sign follows by using two independent draws and averaging
$(x-y)(V'(x)-V'(y))/2$. Hence the consecutive ratios of $M(p)/\Gamma(p)$
strictly decrease as $p$ increases by one. Since
\[
 \gamma_n=\frac{n!}{(2n)!}M(n+1/2)
 =\frac{\sqrt\pi}{4^n}\frac{M(n+1/2)}{\Gamma(n+1/2)},
\]
we have $\gamma_{n+1}^2>\gamma_n\gamma_{n+2}$. All coefficients are
positive. This proves the quadratic assertion.\hfill$\square$

The same calculation applies to every positive eta product with
exponents paired under $d\leftrightarrow N/d$, on centering its
Mellin kernel at weight divided by two. In general this kernel need not
come from an eigenform. For the conductor-$23$ form, the open critical-line
statement would imply real-rootedness in every Jensen degree:
$w\mapsto\Lambda_f(1/2+\sqrt w)$ has order at most $1/2$, and critical-line
zeros correspond to negative real $w$. Hadamard factorization and Jensen's
criterion then give this implication and its converse. The theorem proves
the degree-two part for all shifts. The kernel argument supplements the
prime-sum reduction. It makes no use of its arithmetic estimate.

\subsection{A counterexample to the general cubic kernel criterion}
A possible next step would apply Theorem~1 of D. K. Dimitrov and
F. R. Lucas, \emph{Higher order Tur\'an inequalities for the Riemann
$\xi$-function}, Proc. AMS 139 (2011), 1013--1022,
\href{https://doi.org/10.1090/S0002-9939-2010-10515-4}{p.~1017}.
Together with their normalization and Lemma~1, the printed theorem implies
that every cubic Jensen polynomial is hyperbolic for an admissible
kernel with $(\log K(\sqrt{x}))''<0$ for $x>0$. The following counterexample
addresses this general assertion. It leaves the individual questions for
the Riemann and conductor-$23$ kernels separate.

\textbf{Counterexample.} Put $K_\delta(u)=\exp(-u^2-\delta u^6)$ with
$\delta=10^{-8}$. Its first cubic Jensen polynomial has a nonreal
conjugate pair. This kernel is positive, even and entire, decreases
strictly for $u>0$, and every derivative is a polynomial times $K_\delta$,
hence $O(e^{-|u|^3})$ at infinity. Moreover
$(\log K_\delta(\sqrt{x}))''=-6\delta x<0$.

\textit{Proof.} Use the coefficients $\gamma_n$ above with this new
kernel and put $c_n=4^n\gamma_n/\sqrt\pi$. If $p=n+1/2$ and
$X$ has gamma density $x^{p-1}e^{-x}/\Gamma(p)$, then
$c_n=\mathbb E e^{-\delta X^3}$. With $(p)_m=p(p+1)\cdots(p+m-1)$,
Taylor's inequalities for $e^{-y}$ on $y\ge0$ give
\[
 1-\delta(p)_3\le c_n\le1-\delta(p)_3+
                   \tfrac12\delta^2(p)_6 .
\]
The rescaled Jensen polynomial is
$P(Z)=c_0+3c_1Z+3c_2 Z^2+c_3 Z^3$. These rescalings preserve
real-rootedness. Since $(n+1/2)_3$ is monic cubic in $n$,
its third forward difference is $6$, and $P(-1)=6\delta+O(\delta^2)$.
In $P(W-1)$ the respective coefficients of $W^3,W^2,W,1$ are
$1+O(\delta),O(\delta),O(\delta),6\delta+O(\delta^2)$. The discriminant
is therefore $-972\delta^2+O(\delta^3)$ as $\delta\downarrow0$.
For the stated rational value the accompanying checker uses the displayed
rational endpoints to certify
\[
                 -975<\operatorname{disc}(P)/\delta^2<-969.
\]
Explicitly it encloses $T_1=c_1^2-c_0c_2$, $T_2=c_2^2-c_1c_3$, and
$U=c_1c_2-c_0c_3$ by subtraction of endpoint products. All three lower
endpoints are positive. The final enclosure follows from
$\operatorname{disc}(P)=27(4T_1T_2-U^2)$. \hfill$\square$

One can locate a difficulty in the mean-value argument immediately
before equation (14) on pp.~1020--1021. There
$A(y,z)=(z^2-y^2)K(y)K(z)[B(y)-B(z)]$ with
$B(y)=K'(y)/(yK(y))$. For $y>0$ and $B'(y)<0$, $A(y,z)$ vanishes
quadratically as $z\to y$. Hence $\partial_y A(y,z)/A(y,z)$ has a
simple pole at $z=y$. Its derivative on either open side cannot justify
the mean-value step for an interval with $z_1<y<z_2$.
An argument for the cubic inequalities of a particular arithmetic
kernel must use additional information. The counterexample is a
Mellin-moment calculation and makes no assertion about zeros of a
particular number-theoretic $L$-function.

The obstruction persists with globally strict fourth-derivative curvature.
Take $\widehat K_\delta(u)=\exp[-u^2-\delta(u^4+u^6)]$. Its log has fourth
derivative $-24\delta-360\delta u^2<0$, as does the log of the symmetric
eta kernel above. In the gamma expectation the perturbing polynomial is now
$X^2+X^3$, so the first-order correction to $c_n$ is
$-\delta[(p)_2+(p)_3]$. Its third forward difference is again $-6\delta$.
For every fixed shift $n$ the same translated-cubic calculation gives
discriminant $-972\delta^2+O_n(\delta^3)$ in the $c_n$ normalization.
A single sufficiently small $\delta>0$ therefore violates the cubic
inequalities at any prescribed finite set of shifts. This distinguishes
the information provided by the eta curvature estimate from what a
degree-three argument requires.
